%% ============================================================
%%  SECCIÓN 12 — Observaciones, Sugerencias y Conclusiones
%% ============================================================

\section{Observaciones, Sugerencias y Conclusiones}
\label{sec:conclusiones}

\subsection{Observaciones}
\begin{enumerate}
    \item \textbf{Incertidumbre de Apreciación Visual ("Al Ojo"):}
    La estimación de divisiones en la retícula analógica/digital presenta una limitación intrínseca de resolución humana de aproximadamente $\pm 0{,}1\,\text{div}$. Cuando se opera en escalas grandes (\qty{5}{\volt\per\division}), este error de lectura se magnifica hasta representar un \qty{33}{\percent} de señales pequeñas (\qty{1,5}{\volt}), mientras que en escalas óptimas (\qty{0,2}{\volt\per\division}) su impacto decrece a menos del \qty{1,4}{\percent}.
    
    \item \textbf{Rizado de Rectificación en Fuentes Lineales:}
    Al conmutar el acoplamiento a modo $\text{AC}$ y aumentar la sensibilidad a \qty{5}{\milli\volt\per\division}, se constató cualitativamente la presencia de un rizado de baja frecuencia a \qty{100}{\hertz}/\qty{120}{\hertz}, originado por la rectificación de onda completa de la red eléctrica que no es suprimida en su totalidad por los condensadores de filtrado.

    \item \textbf{Efecto de Carga e Impedancia Interna de Fuentes Reguladas:}
    Las diferentes salidas fijas de la fuente continua entregaron voltajes en circuito abierto sistemáticamente superiores a los valores nominales serigrafiados (\qty{5,220}{\volt} frente a \qty{4,50}{\volt} nominal, \qty{8,030}{\volt} frente a \qty{7,50}{\volt} y \qty{11,150}{\volt} frente a \qty{10,50}{\volt}). Una explicación física plausible para esta elevación de potencial en vacío es el efecto de filtrado capacitivo en reguladores lineales no cargados ($I = 0$, donde $V_{\text{out}} \approx V_{\text{th}}$). No obstante, al no haberse ensayado el circuito con diversas cargas resistivas, esta causa permanece como una hipótesis tentativa bien fundamentada y no como un hecho verificado experimentalmente. Por ello, el contraste y la calibración del osciloscopio se fundamentaron rigurosamente en la medición simultánea del multímetro patrón ($V_{\text{mult}}$), frente al cual el error de lectura en escala óptima fue inferior al \qty{0,45}{\percent}.

    \item \textbf{Fidelidad del Generador DDS:}
    El generador OWON AG 051F demostró que la frecuencia medida en retícula ($f = \qty{250,0 \pm 3,1}{\hertz}$) es plenamente compatible con el valor nominal de \qty{250,00}{\hertz}, confirmando la alta estabilidad temporal del oscilador de cristal de cuarzo integrado ($\pm\qty{1}{\ppm}$).

    \item \textbf{Almacenamiento Digital y Trazabilidad en USB:}
    El registro cuantitativo de los parámetros de retícula, factores de escala y formas de onda en memoria masiva USB (`USB Flash Drive') permitió preservar íntegramente la trazabilidad de los datos experimentales para su posterior procesamiento analítico y reconstrucción gráfica de alta fidelidad científica a 300 DPI, garantizando rigor metodológico y reproducibilidad.
\end{enumerate}

\subsection{Sugerencias}
\begin{enumerate}
    \item \textbf{Uso Sistemático de Cursores Digitales de Pantalla:}
    Para minimizar el error de paralaje y la apreciación visual sobre la retícula, se recomienda emplear la función de cursores de tensión ($\Delta V$) y tiempo ($\Delta t$), la cual aprovecha la resolución completa del ADC ($8$ a $16\text{ bits}$).
    
    \item \textbf{Optimización de la Base de Tiempo Horizontal:}
    Para la medición precisa de períodos y frecuencias, conviene expandir la base de tiempo de modo que un único ciclo abarque entre $8$ y $10\,\text{divisiones}$ horizontales (como se comprobó al pasar de \qty{5}{\milli\second\per\division} a \qty{2}{\milli\second\per\division} para la red de \qty{60}{\hertz}).

    \item \textbf{Adaptación de Impedancias en Generadores de RF:}
    Al conectar generadores de funciones con salida calibrada en $50\,\Omega$, debe tenerse en cuenta que el osciloscopio en modo $\text{High-Z}$ ($1\,\text{M}\Omega$) duplica la amplitud leída en pantalla respecto a la indicada en el dial del generador cuando este asume terminación adaptada.

    \item \textbf{Modo de Promediado de Formas de Onda (\textit{Averaging}):}
    En señales con ruido de fondo o rizado parásito, se sugiere activar el promediado de $16$ o $64$ adquisiciones consecutivas para mejorar la relación señal-ruido ($\text{SNR}$).
\end{enumerate}

\subsection{Conclusiones}
\begin{enumerate}
    \item Se comprobaron los principios de operación y calibración del osciloscopio digital Rohde \& Schwarz HMO1202, verificando que la resolución visual efectiva en la retícula se optimiza maximizando la ocupación vertical de la señal mediante el desplazamiento de la línea base ($\text{GND}$) y la elección de la menor escala de voltios por división sin saturación.

    \item En la medición de tensiones continuas, la calibración vertical demostró un comportamiento lineal en todo el rango evaluado (\qtyrange{1,4}{11,2}{\volt}), con una pendiente de regresión por mínimos cuadrados $m = 1{,}0036 \pm 0{,}0039$ (error estándar $s_m = 0{,}0039$), intercepto $b = -0{,}021 \pm 0{,}026\,\unit{\volt}$ (error estándar $s_b = \qty{0,026}{\volt}$) estadísticamente compatible con cero ($|b|/s_b < 1$) y coeficiente $R^2 = 0{,}99995$, con errores relativos inferiores al \qty{0,45}{\percent} en escalas óptimas.

    \item Para la señal alterna del transformador reductor de \qty{6}{\volt} a \qty{60}{\hertz}, se determinó un voltaje pico a pico de $V_{pp} = \qty{19,50 \pm 0,50}{\volt}$ ($e_r = \qty{0,93}{\percent}$ frente al teórico de \qty{19,32}{\volt}) y un voltaje eficaz calculado de $V_{\text{rms}} = \qty{6,79 \pm 0,14}{\volt}$ ($e_r = \qty{0,61}{\percent}$ frente al multímetro True-RMS). El período medido fue $T = \qty{16,60 \pm 0,20}{\milli\second}$ ($f = \qty{60,24 \pm 0,73}{\hertz}$, $e_r = \qty{0,40}{\percent}$).

    \item Para señales periódicas sintéticas del generador OWON AG 051F a \qty{250}{\hertz}, se validaron experimentalmente las relaciones analíticas entre amplitud pico ($V_p = \qty{4,80 \pm 0,20}{\volt}$) y valor eficaz para onda senoidal ($V_{\text{rms}} = \qty{3,42}{\volt}$), triangular ($V_{\text{rms}} = \qty{2,79}{\volt}$) y cuadrada ($V_{\text{rms}} = \qty{4,80}{\volt}$), con discrepancias inferiores al \qty{0,8}{\percent} frente a la lectura True-RMS.

    \item En el modo graficador $X$-$Y$, se verificó el principio de las figuras de Lissajous para relaciones armónicas enteras $1:1$, $1:2$, $1:3$ y $1:4$, comprobando que la inversión del canal vertical ($\text{CH2 INVERT}$) produce una reflexión especular respecto al eje horizontal $(x, y) \to (x, -y)$, y que potenciales continuos constantes colapsan en puntos estáticos coordenados en el plano cartesiano.
\end{enumerate}
